\subsection{Common Divisor Formulation}

\begin{definition}[Webster Method]
  Let $p_1, \ldots, p_n$ be state populations, $H$ the House size,
  and $N = \sum_i p_i$.
  The Webster apportionment finds a common divisor $d > 0$ such that:
  \[
    \sum_{i=1}^{n} \max\!\left(1, \left\lfloor \frac{p_i}{d} + 0.5 \right\rfloor\right) = H,
  \]
  and assigns each state $s_i = \max(1, \lfloor p_i/d + 0.5 \rfloor)$ seats.
\end{definition}

\noindent
The divisor $d$ is found by binary search or by starting from the natural
divisor $d_0 = N/H$ and adjusting iteratively.
For 2020 Census data, $d_0 = 331{,}108{,}434 / 435 \approx 761{,}169$,
and the Webster divisor is within a small percentage of $d_0$.

\subsection{Priority Queue Formulation}

Equivalently, Webster is a divisor method with priority values:
\[
  P_i^{\text{Web}}(s) = \frac{p_i}{s + 0.5},
\]
giving state $i$ its $(s+1)$-th seat when $p_i/(s+0.5)$ is the
maximum priority.
The first seat priority $P_i^{\text{Web}}(0) = 2p_i$ differs from
Huntington-Hill's first seat priority $P_i^{\text{HH}}(0) = p_i/\sqrt{1\cdot0}$
(which is technically infinite --- every state is guaranteed its first
seat by the constitutional floor and not by the priority queue).
This formulation confirms that Webster is a standard divisor method
and inherits divisor-method paradox immunity.

\subsection{Equivalence to Sainte-Lagu\"e}

In the context of proportional representation for party-list elections,
the Sainte-Lagu\"e method assigns seats to parties using the priority:
\[
  P_j^{\text{SL}}(s) = \frac{V_j}{2s+1},
\]
where $V_j$ is party $j$'s vote total and $s$ its current seat count.
Note that $P_j^{\text{SL}}(s) = V_j / (2s+1)$, and the first-seat
priority is $V_j / 1 = V_j$.

This is proportional to $P_i^{\text{Web}}(s) = p_i/(s+0.5)$:
\[
  P_i^{\text{Web}}(s) = \frac{p_i}{s + 0.5} = \frac{2 p_i}{2s + 1}
  \propto P_j^{\text{SL}}(s).
\]
Since both methods use the same priority order (with votes replacing
populations and parties replacing states), they produce identical
seat allocations whenever the inputs produce the same priority sequence.
The formal equivalence is exact: Webster apportionment equals
Sainte-Lagu\"e apportionment.
This equivalence is documented in \citet{lijphart1994electoral}
and is used in comparative electoral systems analysis.

\subsection{Rounding Threshold Comparison}

For state $i$ with exact quota $q_i = p_i/d$, Webster assigns
$s_i = \lfloor q_i + 0.5 \rfloor$ seats (before applying the constitutional
floor).
The rounding threshold is $s + 0.5$ for each integer $s \geq 0$.
Comparing to Huntington-Hill:
\begin{itemize}
  \item HH threshold: $\sqrt{s(s+1)} < s+0.5$ for all $s \geq 1$
  \item Webster threshold: $s + 0.5$
  \item Gap: $(s+0.5) - \sqrt{s(s+1)} \approx 1/(8s)$ for large $s$
\end{itemize}
A state whose exact quota falls in $({\sqrt{s(s+1)}}, s+0.5)$
receives $s+1$ seats under HH but only $s$ seats under Webster.
This interval is narrow for large $s$ (large states) and wider
for small $s$ (small states), making HH slightly more generous
to small states.
